\section{Sujet n\textsuperscript{o}\,18}
\noindent\begin{minipage}{0.5\linewidth}\footnotesize Office du Baccalauréat du Cameroun\end{minipage}%
\hfill\begin{minipage}{0.45\linewidth}\footnotesize\raggedleft Examen : \textbf{Baccalauréat} · Série : \textbf{C/E}\\
Épreuve : \textbf{Mathématiques} · Durée : \textbf{4\,h} · Coef. : \textbf{7}\end{minipage}
\par\smallskip{\color{onbo}\rule{\linewidth}{1pt}}

\subsection*{Partie A — Évaluation des ressources \hfill (15 points)}

\noindent\textbf{Exercice 1.} \hfill\textit{(3 points)}\\
Au parc national de Waza, on étudie l'observation des grands mammifères lors des
sorties en brousse pour le tourisme écologique. Une sortie permet d'observer un
troupeau d'éléphants avec la probabilité $0{,}6$. Lorsqu'un troupeau d'éléphants est
observé, la probabilité d'apercevoir aussi des lions est $0{,}3$ ; sinon (aucun
éléphant observé), cette probabilité n'est plus que $0{,}5$. On note $E$ l'événement
« observer des éléphants » et $L$ l'événement « observer des lions ».
\begin{enumerate}[label=\arabic*)]
\item Construire l'arbre pondéré décrivant cette situation. \hfill\textit{(0,75 pt)}
\item Calculer la probabilité d'observer à la fois des éléphants et des lions, puis
montrer que $P(L)=0{,}38$. \hfill\textit{(1,25 pt)}
\item Un touriste a observé des lions. Quelle est la probabilité qu'il ait aussi
observé des éléphants ? \hfill\textit{(1 pt)}
\end{enumerate}

\smallskip
\noindent\textbf{Exercice 2.} \hfill\textit{(3 points)}\\
Un éco-gardien doit régler le solde d'un fournisseur d'un montant de
\SI{43}{000} francs CFA à l'aide de billets de \SI{5}{000} francs et de pièces de
\SI{500} francs uniquement. On note $x$ le nombre de billets et $y$ le nombre de
pièces.
\begin{enumerate}[label=\arabic*)]
\item Montrer que le problème revient à résoudre dans $\N^2$ l'équation
$10x+y=86$. \hfill\textit{(0,5 pt)}
\item Résoudre dans $\Z^2$ l'équation $10x+y=86$, puis donner toutes les solutions
$(x\,;\,y)$ dans $\N^2$. \hfill\textit{(1,5 pt)}
\item Le gardien dispose d'au plus $6$ billets de \SI{5}{000} francs. Déterminer
le nombre de billets et de pièces utilisés sachant qu'il veut, de plus, le moins
de pièces possible. \hfill\textit{(1 pt)}
\end{enumerate}

\smallskip
\noindent\textbf{Exercice 3.} \hfill\textit{(4 points)}\\
On modélise la densité de végétation (en unités arbitraires) le long d'un transect
écologique par la fonction $f$ définie sur $\R$ par
$f(x)=x^3-3x-1$, de courbe $(\mathcal{C})$.
\begin{enumerate}[label=\arabic*)]
\item Déterminer les limites de $f$ en $-\infty$ et $+\infty$. \hfill\textit{(0,5 pt)}
\item Calculer $f'(x)$, étudier son signe et dresser le tableau de variations de
$f$. \hfill\textit{(1,5 pt)}
\item Montrer que l'équation $f(x)=0$ admet une unique solution $\alpha$ sur
l'intervalle $[1\,;\,2]$. \hfill\textit{(1 pt)}
\item En déduire un encadrement de $\alpha$ d'amplitude $0{,}1$. \hfill\textit{(1 pt)}
\end{enumerate}

\smallskip
\noindent\textbf{Exercice 4.} \hfill\textit{(5 points)}\\
Dans le plan, trois points de relevé topographique $A$, $B$ et $C$ forment un
triangle tel que $AB=6$, $AC=4$ et $\widehat{BAC}=60^\circ$. Soit $G$ le barycentre
du système $\{(A,2),(B,1),(C,1)\}$.
\begin{enumerate}[label=\arabic*)]
\item Calculer le produit scalaire $\overrightarrow{AB}\cdot\overrightarrow{AC}$,
puis la longueur $BC$. \hfill\textit{(1,25 pt)}
\item Construire le point $G$ et exprimer $\overrightarrow{AG}$ en fonction de
$\overrightarrow{AB}$ et $\overrightarrow{AC}$. \hfill\textit{(1,25 pt)}
\item Pour tout point $M$ du plan, montrer que
$2MA^2+MB^2+MC^2=4MG^2+2GA^2+GB^2+GC^2$. \hfill\textit{(1 pt)}
\item Déterminer et caractériser l'ensemble $(\Gamma)$ des points $M$ du plan
vérifiant $2MA^2+MB^2+MC^2=88$. \hfill\textit{(1,5 pt)}
\end{enumerate}

\subsection*{Partie B — Évaluation des compétences \hfill (5 points)}

\noindent\textit{La conservatrice du parc national de Waza te sollicite, en tant
qu'élève de Terminale~C, pour préparer la saison touristique. Elle veut d'abord
aménager un enclos rectangulaire d'observation \emph{adossé à une rivière} (le côté
longeant la rivière n'est pas clôturé) avec \SI{240}{\metre} de barrière. La
population de gazelles du secteur, estimée à $400$ têtes cette année, augmente
chaque année de \SI{6}{\percent} mais $30$ individus migrent hors du secteur en fin
d'année. Enfin, sur les pistes du parc, \SI{12}{\percent} des véhicules de safari
tombent en panne au cours d'une sortie ; une matinée, $5$ véhicules partent en
sortie de façon indépendante.}

\begin{enumerate}[label=\textbf{Tâche \arabic*.},leftmargin=2.4em]
\item Déterminer les dimensions de l'enclos rectangulaire d'\emph{aire maximale} et
préciser cette aire. \hfill\textit{(1,5 pt)}
\item Déterminer la population de gazelles au bout de $3$ ans et étudier vers
quelle valeur elle tend à long terme. \hfill\textit{(1,5 pt)}
\item Déterminer la probabilité qu'\emph{au moins un} des $5$ véhicules tombe en
panne au cours de la matinée. \hfill\textit{(1,5 pt)}
\end{enumerate}
\par\smallskip{\footnotesize\itshape Présentation générale : 0,5 point.}

% ════════════════════════════════════════════════════
\corrige{du sujet 18}

\noindent\textbf{Partie A — Exercice 1.}
1) Arbre : racine $\to$ $E$ ($0{,}6$) et $\bar E$ ($0{,}4$) ; depuis $E$ : $L$
($0{,}3$), $\bar L$ ($0{,}7$) ; depuis $\bar E$ : $L$ ($0{,}5$), $\bar L$
($0{,}5$).
2) $P(E\cap L)=0{,}6\times0{,}3=0{,}18$. Par la formule des probabilités totales,
$P(L)=P(E\cap L)+P(\bar E\cap L)=0{,}18+0{,}4\times0{,}5=0{,}18+0{,}20=\mathbf{0{,}38}$.
3) $P_L(E)=\dfrac{P(E\cap L)}{P(L)}=\dfrac{0{,}18}{0{,}38}=\dfrac{9}{19}\approx0{,}47$.

\noindent\textbf{Exercice 2.}
1) Le montant payé est $5000x+500y=43000$ francs ; en divisant par $500$ :
$10x+y=86$.
2) $y=86-10x$. Pour tout $x\in\Z$, $(x\,;\,86-10x)$ est solution dans $\Z^2$.
Dans $\N^2$ : il faut $x\ge0$ et $86-10x\ge0$, soit $0\le x\le8$. Les solutions
sont $(x\,;\,86-10x)$ pour $x\in\{0,1,2,\dots,8\}$.
3) $x\le6$ et $y=86-10x$ minimal $\Rightarrow$ prendre $x$ le plus grand possible,
donc $x=6$, $y=86-60=26$. Il utilise \textbf{$6$ billets} de \SI{5}{000} francs et
\textbf{$26$ pièces} de \SI{500} francs.

\noindent\textbf{Exercice 3.}
1) $\lim_{x\to-\infty}f(x)=-\infty$ et $\lim_{x\to+\infty}f(x)=+\infty$ (terme
dominant $x^3$).
2) $f'(x)=3x^2-3=3(x-1)(x+1)$. $f'(x)>0$ sur $]-\infty,-1[\cup]1,+\infty[$ et
$f'(x)<0$ sur $]-1,1[$. Donc $f$ croît sur $]-\infty,-1]$, décroît sur $[-1,1]$,
croît sur $[1,+\infty[$ ; maximum local $f(-1)=-1+3-1=1$, minimum local
$f(1)=1-3-1=-3$.
3) $f$ est continue et strictement croissante sur $[1\,;\,2]$ ;
$f(1)=-3<0$ et $f(2)=8-6-1=1>0$. D'après le théorème des valeurs intermédiaires
(corollaire), $f(x)=0$ admet une unique solution $\alpha\in\,]1\,;\,2[$.
4) $f(1{,}8)=5{,}832-5{,}4-1=-0{,}568<0$ et $f(1{,}9)=6{,}859-5{,}7-1=0{,}159>0$.
Donc $\boxed{1{,}8<\alpha<1{,}9}$ (amplitude $0{,}1$).

\noindent\textbf{Exercice 4.}
1) $\overrightarrow{AB}\cdot\overrightarrow{AC}=AB\cdot AC\cos60^\circ=6\times4\times
\frac12=12$. Avec Al-Kashi : $BC^2=AB^2+AC^2-2\,\overrightarrow{AB}\cdot
\overrightarrow{AC}=36+16-24=28$, donc $BC=2\sqrt7$.
2) $G$ barycentre de $\{(A,2),(B,1),(C,1)\}$, masse totale $4$. De
$2\overrightarrow{GA}+\overrightarrow{GB}+\overrightarrow{GC}=\vec0$ on tire
$\overrightarrow{AG}=\frac14(2\overrightarrow{AA}+\overrightarrow{AB}+
\overrightarrow{AC})=\frac14(\overrightarrow{AB}+\overrightarrow{AC})$. ($G$ est le
milieu du segment joignant $A$ au milieu de $[BC]$.)
3) Pour tout $M$ : $2MA^2+MB^2+MC^2=\sum\alpha_i\overrightarrow{MG_i}^2$. En écrivant
$\overrightarrow{MA}=\overrightarrow{MG}+\overrightarrow{GA}$ etc. et en utilisant
$2\overrightarrow{GA}+\overrightarrow{GB}+\overrightarrow{GC}=\vec0$, le double
produit s'annule, d'où
$2MA^2+MB^2+MC^2=(2+1+1)MG^2+2GA^2+GB^2+GC^2=4MG^2+2GA^2+GB^2+GC^2$.
4) Calcul de la constante $k=2GA^2+GB^2+GC^2$. Avec
$\overrightarrow{AG}=\frac14(\overrightarrow{AB}+\overrightarrow{AC})$ :
$GA^2=\frac1{16}\big(AB^2+AC^2+2\,\overrightarrow{AB}\cdot\overrightarrow{AC}\big)
=\frac1{16}(36+16+24)=\frac{76}{16}=\frac{19}{4}$.
Or $\overrightarrow{GB}=\overrightarrow{GA}+\overrightarrow{AB}$ et
$\overrightarrow{GC}=\overrightarrow{GA}+\overrightarrow{AC}$ ; on calcule
$GB^2=GA^2+AB^2+2\overrightarrow{GA}\cdot\overrightarrow{AB}$ et
$GC^2=GA^2+AC^2+2\overrightarrow{GA}\cdot\overrightarrow{AC}$, avec
$\overrightarrow{GA}=-\frac14(\overrightarrow{AB}+\overrightarrow{AC})$ :
$\overrightarrow{GA}\cdot\overrightarrow{AB}=-\frac14(AB^2+\overrightarrow{AB}\cdot
\overrightarrow{AC})=-\frac14(36+12)=-12$ ;
$\overrightarrow{GA}\cdot\overrightarrow{AC}=-\frac14(\overrightarrow{AB}\cdot
\overrightarrow{AC}+AC^2)=-\frac14(12+16)=-7$.
Donc $GB^2=\frac{19}4+36-24=\frac{19}4+12=\frac{67}4$ et
$GC^2=\frac{19}4+16-14=\frac{19}4+2=\frac{27}4$.
Ainsi $k=2\times\frac{19}4+\frac{67}4+\frac{27}4=\frac{38}4+\frac{67}4+\frac{27}4
=\frac{132}{4}=33$.
L'équation $4MG^2+33=88$ donne $MG^2=\frac{55}{4}$, soit $MG=\frac{\sqrt{55}}2$.
$(\Gamma)$ est le \textbf{cercle de centre $G$ et de rayon $\frac{\sqrt{55}}2\approx
3{,}71$}.

\smallskip
\noindent\textbf{Partie B — Situation.}
\emph{Tâche 1.} Largeur $x$ (les deux côtés perpendiculaires à la rivière),
longueur $y=240-2x$ (un seul côté parallèle à la rivière, l'autre étant la rivière).
Aire $A(x)=x(240-2x)=240x-2x^2$ ; $A'(x)=240-4x=0$ en $x=60$. Dimensions :
$x=\SI{60}{\metre}$, $y=\SI{120}{\metre}$ ; aire maximale
$60\times120=\SI{7200}{\metre\squared}$.

\emph{Tâche 2.} Population de l'année $n$ : $u_{n+1}=1{,}06\,u_n-30$, $u_0=400$.
Point fixe $\ell=1{,}06\ell-30\Rightarrow\ell=500$. Posons $v_n=u_n-500$ : c'est une
suite géométrique de raison $1{,}06$, $v_0=-100$, donc $u_n=500-100\times1{,}06^{\,n}$.
$u_3=500-100\times1{,}06^3=500-100\times1{,}191016=500-119{,}1\approx\textbf{380}$
gazelles. Comme $1{,}06^{\,n}\to+\infty$, $u_n\to-\infty$ : le modèle prévoit la
\textbf{disparition} de la population (elle décroît, $u_n<u_0$, et finit par
s'annuler ; concrètement la population s'éteint).

\emph{Tâche 3.} Le nombre $X$ de véhicules en panne suit la loi binomiale
$\mathcal{B}(5\,;\,0{,}12)$. $P(\text{au moins une panne})=1-P(X=0)=1-(0{,}88)^5
\approx1-0{,}5277=\mathbf{0{,}472}$, soit environ \SI{47,2}{\percent}.
